\documentclass[aps,prb,reprint,amsmath,amssymb,nofootinbib]{revtex4-2}
\usepackage{bm}
\usepackage{booktabs}
\usepackage[colorlinks=true,linkcolor=blue,citecolor=blue,urlcolor=blue]{hyperref}
\newcommand{\ket}[1]{\lvert#1\rangle}
\newcommand{\bra}[1]{\langle#1\rvert}
\newcommand{\Diss}[1]{\mathcal{D}[#1]}
\newcommand{\tr}{\operatorname{Tr}}
% Generated from config/noise_model.json; do not edit.
\newcommand{\ChiLA}{1}
\newcommand{\ChiLB}{1}
\newcommand{\ChiBA}{1}
\newcommand{\ChiBB}{1.3}
\newcommand{\DeltaBus}{8.05}
\newcommand{\CavBound}{0.25}
\newcommand{\QBound}{1}
\newcommand{\CavityLife}{600}
\newcommand{\QOne}{50}
\newcommand{\QTwo}{30}
\newcommand{\Cutoff}{24}
\newcommand{\JumpStep}{0.005}

\begin{document}
\title{Methods: a driven--dissipative model for the Nov17 two-cavity circuit}
\date{November 17, 2026}
\begin{abstract}
We specify the system Hamiltonian and Markovian noise model used for a
two-cavity circuit controlled by two local transmons and one shared bus
transmon. Each cavity encodes three binary spins, giving six encoded spins
and three auxiliary physical qubits. The specification fixes operator,
rotating-frame, drive-amplitude, and dephasing conventions so that
density-matrix and quantum-trajectory calculations use the same generator.
Numerical parameters are an illustrative baseline; only the cavity lifetime
was supplied as a device value. This methods note is a model specification,
not an experimental characterization or a published journal article.
\end{abstract}
\maketitle

\section{Methods}
\subsection{Architecture, Hilbert space, and encoding}

The device contains two bosonic storage modes, labeled $A$ and $B$, local
control transmons $q_A$ and $q_B$, and a transmon $q_b$ dispersively coupled
to both modes. Figure~\ref{fig:architecture} shows the retained coupling
graph. We approximate each transmon by $\{\ket{g},\ket{e}\}$ and write
$a_j$ for the cavity annihilation operator and $n_j=a_j^\dagger a_j$ for
photon number. We use the quantum-optics Pauli convention
\begin{align}
 \sigma_r^z&=\ket{e_r}\bra{e_r}-\ket{g_r}\bra{g_r},\\
 \sigma_r^-&=\ket{g_r}\bra{e_r},\qquad
 \sigma_r^+=\ket{e_r}\bra{g_r},\\
 \sigma_r^x&=\sigma_r^++\sigma_r^-,\qquad
 \sigma_r^y=-i(\sigma_r^+-\sigma_r^-).
\end{align}
Thus $\sigma_r^z\ket g=-\ket g$, $\sigma_r^z\ket e=\ket e$, and
$\sigma_r^\pm=(\sigma_r^x\pm i\sigma_r^y)/2$.
The physical-transmon operators $\sigma_r^\alpha$ are distinct from the
encoded-spin operators $X_i,Z_i$. Tensor-product identities are implicit.
All physical transmons are initialized in $g$.

\begin{figure}[t]
\centering
\setlength{\unitlength}{1pt}
\begin{picture}(235,104)
\put(34,76){\oval(57,25)}\put(13,73){$q_A$ (local)}
\put(199,76){\oval(57,25)}\put(179,73){$q_B$ (local)}
\put(34,24){\oval(63,27)}\put(12,21){Cavity $A$}
\put(199,24){\oval(63,27)}\put(178,21){Cavity $B$}
\put(116,24){\oval(49,27)}\put(100,21){Bus $q_b$}
\put(34,38){\line(0,1){25}}\put(42,48){$\chi_A$}
\put(199,38){\line(0,1){25}}\put(207,48){$\chi_B$}
\put(66,24){\line(1,0){25}}\put(69,35){$\chi_{bA}$}
\put(141,24){\line(1,0){26}}\put(143,35){$\chi_{bB}$}
\end{picture}
\caption{Retained dispersive couplings. The two cavities encode six logical
spins in total. The three transmons are auxiliary physical qubits; the bus
is shared, whereas each local transmon addresses one cavity. No direct
cavity--cavity exchange or cross-Kerr is included in the baseline.}
\label{fig:architecture}
\end{figure}

For cavity cutoffs $N_A$ and $N_B$, the numerical tensor ordering is
\begin{align}
 \mathcal H={}&\mathbb C_{q_b}^{2}\otimes\mathbb C_{q_A}^{2}
 \otimes\mathbb C_{c_A}^{N_A}\otimes\mathbb C_{q_B}^{2}
 \otimes\mathbb C_{c_B}^{N_B},\\
 d={}&8N_A N_B.
\end{align}
The logical subspace of each cavity uses Fock levels $0$--$7$, with
$\ket{b_2b_1b_0}\leftrightarrow\ket{4b_2+2b_1+b_0}$.
The spin groups are $\{0,2,4\}$ and $\{1,3,5\}$, in local
least-significant-to-most-significant-bit order. The encoded space thus has
dimension $8^2=64$. Simulations allow excursions beyond it; projection at
analysis time is not accompanied by renormalization or postselection.

\subsection{Dispersive Hamiltonian and control fields}

Dispersive cavity--transmon descriptions and microwave control of bosonic
modes are standard in circuit QED.\cite{Blais2021} A particularly relevant
experimental precedent is the two-cavity/shared-transmon Hamiltonian in
Eq.~(3) of Wang \emph{et al.},\cite{Wang2016} which also retains a third
transmon level. We retain only its $g/e$ structure and add a local
dispersive control transmon for each cavity. This is a modeling extension;
the cited experiment does not implement our full five-component graph or
the Nov17 encoded-spin ansatz.

Let $h=H_{\mathrm{sys}}/\hbar$ denote an angular-frequency Hamiltonian.
After transforming to the cavity and transmon drive frames and making a
rotating-wave approximation, we use
\begin{equation}
 \boxed{\begin{aligned}
 h(t)={}&\sum_{j=A,B}\widetilde\Delta_{c,j}n_j
       +\sum_{r=A,B,b}\frac{\Delta_{q,r}}2\sigma_r^z\\
 &-\frac12\sum_{r,j}\chi_{rj}n_j\sigma_r^z\\
 &+\sum_{j=A,B}[\epsilon_j a_j^\dagger+\epsilon_j^*a_j]\\
 &+\sum_{r=A,B,b}[\Omega_r\sigma_r^++\Omega_r^*\sigma_r^-]
       +h_{\mathrm{res}} .
 \end{aligned}}
 \label{eq:system}
\end{equation}
This is the total effective Hamiltonian for all three transmons and both
cavities; the envelopes depend on time. Its coupling matrix (rows
$q_A,q_B,q_b$; columns $A,B$) is
\begin{equation}
 (\chi_{rj})=\begin{pmatrix}
 \chi_A&0\\0&\chi_B\\\chi_{bA}&\chi_{bB}
 \end{pmatrix}.
\end{equation}
We write $h=h_0+h_{\mathrm d}+h_{\mathrm{res}}$ to separate the fixed
couplings from the applied drives.
The baseline drift is explicitly
\begin{align}
 h_0={}&\sum_{j=A,B}\widetilde\Delta_{c,j}n_j
       +\frac12\sum_{r=A,B,b}\Delta_{q,r}\sigma_r^z\nonumber\\
     &-\frac{\chi_A}{2}n_A\sigma_A^z
      -\frac{\chi_B}{2}n_B\sigma_B^z\nonumber\\
     &-\frac12(\chi_{bA}n_A+\chi_{bB}n_B)\sigma_b^z .
 \label{eq:drift}
\end{align}
The detunings in the simulation configuration refer to the cavity
frequencies with all transmons in $g$. In Pauli notation this requires
\begin{equation}
 \widetilde\Delta_{c,j}=\Delta_{c,j}-\frac12\sum_r\chi_{rj}.
 \label{eq:cavityshift}
\end{equation}
Here $\Delta_{c,j}$ and $\Delta_{q,r}$ are the configured mode-minus-drive
angular detunings. We omit only the global scalar
$\frac12\sum_r\Delta_{q,r}$ from $H/\hbar$; this does not affect any
observable or the master equation. Equation~\eqref{eq:cavityshift} must
be retained: a photon-number-dependent frequency shift is not a global
energy offset. The $1/2$ in the coupling makes $\chi_{rj}$ the full
transmon transition shift per photon. This is the usual Pauli dispersive
form, with our explicitly negative shift convention.\cite{Blais2021}

The complex drive envelopes are $\epsilon_j=\epsilon_{jI}+i\epsilon_{jQ}$
and $\Omega_r=\Omega_{rI}+i\Omega_{rQ}$. We write
\begin{align}
 h_{\mathrm d}(t)={}&\sum_{j=A,B}
       [\epsilon_j(t)a_j^\dagger+\epsilon_j^*(t)a_j]\nonumber\\
 &+\sum_{r=A,B,b}
       [\Omega_r(t)\sigma_r^++\Omega_r^*(t)\sigma_r^-].
 \label{eq:drives}
\end{align}
Equivalently, the cavity quadratures multiply $(a+a^\dagger)$ and
$i(a^\dagger-a)$. With $\Omega_r=\Omega_{rI}+i\Omega_{rQ}$, the transmon
drive is $\Omega_{rI}\sigma_r^x-\Omega_{rQ}\sigma_r^y$.
The minus sign follows from the Pauli convention defined above; it
preserves the complex envelopes and I/Q samples of the stored waveforms.
There are ten independent real controls. A resonant constant envelope
has angular Rabi frequency $2|\Omega_r|$, so $\Omega_r$ is half the
complex Rabi-frequency envelope in a convention that writes the drive
with an overall $1/2$.

For later residual-coupling studies, the implementation also permits
\begin{align}
 h_{\mathrm{res}}={}&\sum_{j=A,B}\frac{K_j}{2}n_j(n_j-1)
       +K_{AB}n_A n_B\nonumber\\
 &+J_{AB}(a_A^\dagger a_B+a_A a_B^\dagger).
 \label{eq:residual}
\end{align}
All coefficients in Eq.~\eqref{eq:residual} are zero in the reported
baseline. Their signs follow the equation as written. They are optional
calibration parameters, not measured zeros. In particular, parking the bus
in $g$ makes its dispersive contribution cancel the corresponding
part of Eq.~\eqref{eq:cavityshift}, but does not in general
eliminate higher-order interactions of a real device. Independent local
control is exact only in the ideal retained graph with negligible residuals.

\begin{table}[t]
\caption{Symbols used in the system model. All detunings, dispersive shifts,
and drive coefficients in $h$ are angular frequencies.}
\label{tab:symbols}
\begin{ruledtabular}
\begin{tabular}{p{0.23\columnwidth}p{0.68\columnwidth}}
$a_j,n_j$ & Cavity ladder and number operators, $j=A,B$.\\
$q_A,q_B,q_b$ & Local control transmons and shared bus.\\
$\sigma_r^z$ & Transmon Pauli operator; $g:-1$, $e:+1$.\\
$\sigma_r^\pm$ & Transmon raising/lowering operators.\\
$\Delta_{c,j},\Delta_{q,r}$ & Configured drive-frame detunings.\\
$\widetilde\Delta_{c,j}$ & Pauli-form cavity detuning, Eq.~\eqref{eq:cavityshift}.\\
$\chi_A,\chi_B$ & Local dispersive shifts.\\
$\chi_{bA},\chi_{bB}$ & Shared-bus dispersive shifts.\\
$\epsilon_j,\Omega_r$ & Complex cavity and transmon drives.\\
$K_j,K_{AB},J_{AB}$ & Optional self-Kerr, cross-Kerr, and exchange.\\
$\kappa_j,\gamma_{1r},\gamma_{\phi r}$ & Photon-loss, relaxation, and pure-dephasing rates.\\
\end{tabular}
\end{ruledtabular}
\end{table}

\subsection{Numerical baseline and spectral selectivity}

We use time in microseconds. Quoted frequencies in MHz are cycles per
microsecond and are multiplied by $2\pi$ in $h$. In contrast, dissipative
rates $1/T$ are used directly in $\mu\mathrm{s}^{-1}$, without $2\pi$.
The baseline values are collected in Table~\ref{tab:parameters}. For the
configured zero local and cavity detunings, the Pauli-form cavity
coefficients are $\widetilde\Delta_{c,A}/(2\pi)=-1$ MHz and
$\widetilde\Delta_{c,B}/(2\pi)=-1.15$ MHz. Thus
\begin{align}
 \frac{h_0}{2\pi\,\mathrm{MHz}}={}&-n_A-1.15n_B
       -\frac{\ChiLA}{2}n_A\sigma_A^z
       -\frac{\ChiLB}{2}n_B\sigma_B^z\nonumber\\
 &+\frac12(\DeltaBus-\ChiBA\,n_A-\ChiBB\,n_B)\sigma_b^z .
 \label{eq:numerical}
\end{align}
This differs from the stored Hamiltonian matrix only by the omitted
global scalar $\Delta_{q,b}/2$.

The bus detuning is chosen so that its carrier lies at the center of the
logical joint-number manifold. Its conditional transition detuning is
$\Delta_{q,b}-\chi_{bA}n_A-\chi_{bB}n_B$. At the assumed unequal shifts,
the 64 logical pairs are distinct, with minimum separation $0.1$ MHz.
Equal shifts would instead leave many pairs spectrally degenerate; an
undriven bus in either energy eigenstate supplies separated frequency
shifts, rather than an arbitrary joint phase.

The shift values, coherence times, and ideal connectivity used here are
assumptions; the references do not establish them for this device.

\subsection{Which terms implement each operation?}

\paragraph{Joint two-cavity SNAP.}
Heeres \emph{et al.}\cite{Heeres2015} use a dispersive shift to resolve
Fock-number-conditioned transmon transitions and a transmon drive to
imprint selected phases, with the transmon returned to $g$. Their
Eq.~(1) writes the dispersive term with $\ket e\bra e$ and a signed
shift. Substituting $\ket e\bra e=(I+\sigma^z)/2$ gives the same Pauli
structure used here, with the associated cavity-frequency shift and
our chosen sign of $\chi$. Their experiment controls one storage mode
(the other resonator is for readout); our joint phase gate is a two-storage-mode
extension, using the shared dispersive structure of Ref.~\onlinecite{Wang2016}.

During the joint gate only the bus drive is applied:
$\epsilon_A=\epsilon_B=\Omega_A=\Omega_B=0$, $\Omega_b(t)\ne0$.
With both local transmons in $g$, the ideal Hamiltonian decomposes into
independent blocks labeled by $(n,m)$:
\begin{align}
 h_b^{(n,m)}(t)={}&-\frac{\chi_{bA}n+\chi_{bB}m}{2}I
                     +\frac{\delta_{nm}}2\sigma_b^z\nonumber\\
                 &+\Omega_b(t)\sigma_b^++\Omega_b^*(t)\sigma_b^-,\\
 \delta_{nm}={}&\Delta_{q,b}-\chi_{bA}n-\chi_{bB}m.
 \label{eq:busblock}
\end{align}
Here the configured cavity detunings and residuals are zero; constant
spectator-transmon energies have been omitted. Although the first term
is proportional to the identity within one bus block, its coefficient
depends on $(n,m)$ and must be retained when computing cavity phases.
The same
waveform must steer every block to
\begin{equation}
 U_b^{(n,m)}(T)\ket g\simeq e^{-i\Phi_{nm}}\ket g,
\end{equation}
producing the joint SNAP target
\begin{equation}
 S_{AB}(\Phi)=\sum_{n,m=0}^{7}e^{-i\Phi_{nm}}
                  \ket{n,m}\bra{n,m}.
 \label{eq:jointsnap}
\end{equation}
Dispersive coupling provides number selectivity; the driven closed bus
motion creates the phase. Neither free evolution with the bus fixed in
$g$ nor two independent local SNAP gates creates an arbitrary joint
phase. A phase table with genuinely nonseparable entries can entangle
the cavities. Our stored pulses use optimal control of the bus I/Q
envelopes to approximate Eq.~\eqref{eq:jointsnap}; they are not a literal
sequence of ideal spectrally isolated geometric SNAP pulses.
The small $0.1$-MHz minimum splitting makes finite-duration selectivity a
control problem, not a guarantee from unequal shifts alone.

Zhang \emph{et al.}\cite{Zhang2019} (arXiv:1808.07960) treat bichromatically driven, weakly anharmonic
transmons and engineered bilinear mode coupling, including
$g_{\rm BS}a_A^\dagger a_B+g_{\rm BS}^*a_Aa_B^\dagger$.
That operation exchanges photons. Equation~\eqref{eq:jointsnap} is
number-diagonal and uses a different control mechanism. We do not add
that beam-splitter term or identify their two-tone SWAP protocol with
our SNAP. Their analysis also shows why strong-drive dissipation can
require a richer model than constant bare lifetimes.

\paragraph{Local GRAPE for preparation or a Hadamard.}
Park the bus in $g$ and set $\Omega_b=0$. With baseline residuals zero,
the two local controls act in parallel under
\begin{align}
 h_j(t)={}&(\Delta_{c,j}-\chi_j/2)n_j
             +\frac{\Delta_{q,j}}2\sigma_j^z\nonumber\\
 &-\frac{\chi_j}{2}n_j\sigma_j^z\nonumber\\
 &+\epsilon_j(t)a_j^\dagger+\epsilon_j^*(t)a_j\nonumber\\
 &+\Omega_j(t)\sigma_j^++\Omega_j^*(t)\sigma_j^-,
 \quad j=A,B.
 \label{eq:localgrape}
\end{align}
GRAPE optimizes the four real drive quadratures per cavity; it is an
optimization method, not another interaction term. The cavity drive
mixes Fock populations, while the local transmon and its dispersive
coupling supply number-dependent control.\cite{Krastanov2015}
The preparation target is only
\begin{equation}
 U_j(T)\ket{g,0}\simeq\ket g\frac1{\sqrt8}
                         \sum_{n=0}^{7}\ket n.
\end{equation}
The coherent Hadamard target is stronger:
$U_j(T)\ket{g,n}\simeq\ket g H^{\otimes3}\ket n$ for every logical
input $n$. Thus the shorter vacuum-to-plus preparation pulse need not
implement a general Hadamard. During local diagonal-shell SNAP, set
$\epsilon_j=0$ and shape $\Omega_j$ using the same local dispersive term.
The drift couplings remain present in all stages; the schedule changes
which drives are applied. Full noisy replay retains the whole
Hamiltonian even if an auxiliary leaves $g$.

\begin{table}[t]
\caption{Current simulation baseline. Only the cavity lifetime was
user-supplied. All other values are illustrative, not fitted device data.}
\label{tab:parameters}
\begin{ruledtabular}
\begin{tabular}{ll}
Parameter & Value\\ \hline
$\chi_A/(2\pi),\chi_B/(2\pi)$ & $\ChiLA,\ChiLB$ MHz\\
$\chi_{bA}/(2\pi),\chi_{bB}/(2\pi)$ & $\ChiBA,\ChiBB$ MHz\\
$\Delta_{q,b}/(2\pi)$ & $\DeltaBus$ MHz\\
Other configured detunings & $0$\\
$|\epsilon_{jI,Q}|/(2\pi)$ bound & $\CavBound$ MHz\\
$|\Omega_{rI,Q}|/(2\pi)$ bound & $\QBound$ MHz\\
Cavity $T_1$ & $\CavityLife\,\mu\mathrm{s}$\\
Each transmon $T_1,T_2$ & $\QOne,\QTwo\,\mu\mathrm{s}$\\
Kerr, exchange, thermal population & $0$ in baseline\\
Default cavity cutoff & $N_A=N_B=\Cutoff$\\
Full-circuit jump step bound & $\JumpStep\,\mu\mathrm{s}$\\
\end{tabular}
\end{ruledtabular}
\end{table}

\subsection{Lindblad generator and rate conventions}

We assume independent, zero-temperature Markovian baths. The density
operator of both cavities and all three transmons obeys\cite{Lindblad1976}
\begin{align}
 \dot\rho={}&-i[h(t),\rho]
       +\sum_{j=A,B}\kappa_j\Diss{a_j}\rho\nonumber\\
 &+\sum_{r=A,B,b}\left[
       \gamma_{1r}\Diss{\sigma_r^-}\rho
       +\frac{\gamma_{\phi r}}{2}\Diss{\sigma_r^z}\rho
       \right],
 \label{eq:lindblad}\\
 \Diss{L}\rho={}&L\rho L^\dagger-\frac12\{L^\dagger L,\rho\}.
\end{align}
For each transmon, measured energy and coherence lifetimes would give
\begin{equation}
 \gamma_{1r}=\frac1{T_{1r}},\qquad
 \gamma_{\phi r}=\frac1{T_{2r}}-\frac1{2T_{1r}}\geq0.
 \label{eq:rates}
\end{equation}
Thus $T_2$ is not the pure-dephasing time. In the baseline,
\begin{align}
 \kappa_j&=1/600=0.00166667\;\mu\mathrm{s}^{-1},\\
 \gamma_{1r}&=0.0200000\;\mu\mathrm{s}^{-1},\\
 \gamma_{\phi r}&=0.0233333\;\mu\mathrm{s}^{-1},\\
 T_{\phi r}&=42.8571\;\mu\mathrm{s}.
\end{align}
The implementation uses eight collapse operators,
\begin{align}
 L_{c,j}&=\sqrt{\kappa_j}\,a_j,\\
 L_{1,r}&=\sqrt{\gamma_{1r}}\,\sigma_r^-,\\
 L_{\phi,r}&=\sqrt{\gamma_{\phi r}/2}\,\sigma_r^z .
 \label{eq:collapse}
\end{align}
This convention gives the coefficient $\gamma_{\phi r}/2$ in
Eq.~\eqref{eq:lindblad}, so an otherwise idle transmon has coherence
decay $\exp(-t/T_{2r})$. The simulation's excitation-diagonal dephasing
operator generates exactly the same dissipator: adding a real multiple
of the identity to a Hermitian collapse operator does not change its
dissipator. The displayed Pauli convention is therefore physically
equivalent to the retained numerical noise model.

Here cavity $T_1$ is the photon-number energy lifetime, so
$\langle n_j(t)\rangle=e^{-\kappa_jt}\langle n_j(0)\rangle$.
The first photon-loss waiting time in $\ket n$ is $T_1/n$, rather than a
state-independent encoded-spin lifetime. It is $85.7\,\mu\mathrm{s}$ for
$\ket7$ and $42.9\,\mu\mathrm{s}$ for $\ket{7,7}$ at equal cavity
lifetimes. Excursions above the logical levels during control must also
be included when assessing loss exposure.

Cavity number dephasing is zero in the baseline. It can be added as
$\eta_j\Diss{n_j}\rho$, corresponding to
$L_{n,j}=\sqrt{\eta_j}\,n_j$ and off-diagonal Fock decay
$\exp[-\eta_j(n-m)^2t/2]$.
Finite-temperature excitation, correlated baths, higher transmon levels,
drive-dependent noise, and detector/reset errors are not inferred from
the supplied lifetimes and are not included in Eq.~\eqref{eq:lindblad}.

\subsection{Nov17 pulse scheduling and the simulated spin model}

The device Hamiltonian~\eqref{eq:system} is distinct from the encoded
variational objective,
\begin{equation}
 H_{\mathrm{SK}}=\sum_{i<j}J_{ij}Z_iZ_j-B\sum_iX_i.
\end{equation}
The latter is not an additional drift term during microwave pulses.
We now specify the matrices used by the unchanged Nov3 ansatz. In fixed
spin-energy units, $J=J^T$, $J_{ii}=0$, and independent upper-triangle
entries satisfy $J_{ij}\sim\mathcal N(0,1/6)$. Set
$G_A=\{0,2,4\}$ and $G_B=\{1,3,5\}$, and define
\begin{align}
 K^{(0)}_{ij}&=\begin{cases}
       \sum_{s=0}^{5}J_{is}J_{sj},&i\ne j,\\0,&i=j,
       \end{cases}\\
 \bar J_{AB}&=\frac19\sum_{i\in G_A}\sum_{j\in G_B}J_{ij},\\
 M^{(0)}_{ij}&=\begin{cases}
       \bar J_{AB},&i\in G_A,j\in G_B\text{ or vice versa},\\
       0,&\text{otherwise}.
       \end{cases}
\end{align}
For a symmetric matrix $Q$, let
\begin{align}
 r(Q)&=\left[\frac1{15}\sum_{i<j}Q_{ij}^2\right]^{1/2},\\
 K&=\frac{r(J)}{r(K^{(0)})}K^{(0)},\qquad
 M=\frac{r(J)}{r(M^{(0)})}M^{(0)}.
\end{align}
A zero template is left zero. Thus $K$ is a normalized off-diagonal
\emph{matrix square}, not the square of the many-spin Hamiltonian;
$M$ is a normalized constant cross-cavity template. Both still multiply
two-spin Pauli operators. The shared cost gate in every layer is
\begin{align}
 A&=\gamma(J+aK+bM),\\
 C_Z&=\exp\!\left[-i\sum_{i<j}A_{ij}Z_iZ_j\right].
 \label{eq:cost}
\end{align}
The parameters $\gamma,a,b$ are shared across layers, not independently
chosen for each layer. Their combination makes $A_{ij}$ dimensionless
phase angles. The local shell is
\begin{align}
 \hat k_j&=\sum_{i\in G_j}\frac{I-Z_i}{2},\quad j=A,B,\\
 f_\ell(k)&=b_{\ell1}k+b_{\ell2}{k\choose2}
                         +b_{\ell3}{k\choose3},\\
 S_\ell^Z&=\exp\!\left[-i\sum_{j=A,B}f_\ell(\hat k_j)\right].
\end{align}
Here $k$ is the Hamming weight of the three bits, not the Fock number;
the same three $b_{\ell p}$ apply to both cavities. This shell contains
local $Z$, $ZZ$, and $ZZZ$ contributions and must not be replaced by
single-$Z$ rotations alone. With $W=H^{\otimes6}$, the evolving state is
\begin{align}
 \ket{\psi_0}&=W\ket0^{\otimes6}=\ket+^{\otimes6},\\
 \ket{\psi_\ell}&=W S_\ell^Z W C_Z\ket{\psi_{\ell-1}},\\
 \ket{\psi_L(\bm\theta)}&=
  (W S_L^Z W C_Z)\cdots(W S_1^Z W C_Z)W\ket0^{\otimes6}.
 \label{eq:ansatzstate}
\end{align}
Operators act from right to left, with $3+3L$ variational parameters.
Each $W$ is two local three-bit Hadamards in parallel; the first $W$ on
vacuum can use the less restrictive preparation target above.

For implementation, write the bits of $n$ on $G_A$ and those of $m$ on
$G_B$ as $s_i(n,m)\in\{0,1\}$. Then the joint SNAP phase table is
\begin{equation}
 \Phi_{nm}=\sum_{i<j}A_{ij}
           [1-2s_i(n,m)][1-2s_j(n,m)].
\end{equation}
All 15 cost terms commute: six are within cavities (three per cavity)
and nine are cross-cavity. The retained implementation compiles their
\emph{complete} phase table into one bus pulse. They can alternatively
be factored into local and joint SNAPs, but that is not the stored
all-15 pulse schedule.

Omitting the last $W$ changes the output coordinates to the original-$X$
frame. Native Fock readout then supplies original-$X$ observables, whereas
a final calibrated measurement basis change supplies original-$Z$
correlations. This is the same logical state with consistently transformed
observables. A phase-free readout transformation is suitable at the final
measurement boundary, but an interior $W$ requires coherent control.

\subsection{Numerical realization and scope of existing simulations}

The numerical Hamiltonian is piecewise constant in the ten real drive
quadratures, with fixed drift present throughout each operation. For
$N_A=N_B=24$, the Hilbert-space dimension is $4608$. The same Hamiltonian
and collapse operators can be used in a density-matrix ordinary
differential equation, without constructing a $d^2\times d^2$
Liouvillian matrix.

Existing full-circuit noisy runs instead use the quantum-jump
unraveling,\cite{Dalibard1992} whose ensemble approaches
Eq.~\eqref{eq:lindblad}. Between jumps the unnormalized state evolves
under
\begin{equation}
 h_{\mathrm{eff}}(t)=h(t)-\frac{i}{2}\sum_\mu L_\mu^\dagger L_\mu.
\end{equation}
The norm determines the no-jump probability; jump channels are selected
with weights proportional to $\|L_\mu\ket\psi\|^2$.
Reported Nov17 comparisons used 512 trajectories and a maximum jump
step of $5$ ns on GPU~1. No failed trajectories were discarded.
Finite-step error and Monte Carlo uncertainty remain distinct from
uncertainty in the assumed device parameters. Local preparation was
also checked by density-matrix propagation with symmetric unitary/noise
splitting. A full-circuit dense density-matrix integration has not yet
been performed.

Algebraic verification on 20 saved and 120 random parameter sets confirms
that the Nov17 ordering reproduces the original Nov3 ansatz and both
output frames to numerical roundoff. Pulse searches and noisy
preparation comparisons are finite tests, not certificates of optimal
control. Cavity-cutoff convergence tests address numerical truncation;
they do not establish the physical validity of a dispersive model at
high photon number. That requires the actual coupling-to-detuning
ratios, transmon anharmonicities, and critical photon numbers.
The present model is therefore a reproducible starting point for future
Lindblad studies and device calibration, not a substitute for that
calibration.


\subsection{Using the Hamiltonian for noisy evolution}

Insert each compiled waveform into Eq.~\eqref{eq:system} to obtain the
piecewise time-dependent $h(t)$. Initialize the two cavities in vacuum
and all three transmons in $g$, so $\rho(0)=\ket{g,g,0,g,0}
\bra{g,g,0,g,0}$ in our tensor order. Integrate Eq.~\eqref{eq:lindblad}
through preparation, cost, mixer, and any readout pulse with the same
drift and baths active throughout. The baseline noise channels are:
\begin{itemize}
 \item \emph{Cavity photon loss:} energy escapes to the environment;
 each cavity uses $T_{1c}=600\,\mu\mathrm{s}$,
 $L=\sqrt{1/T_{1c}}\,a$. Photon loss changes the encoded Fock state.
 \item \emph{Transmon relaxation:} each local transmon and the bus can
 decay $e\to g$; the assumed $T_{1q}=50\,\mu\mathrm{s}$ sets
 $L=\sqrt{1/T_{1q}}\,\sigma^-$.
 \item \emph{Transmon pure dephasing:} unresolved frequency fluctuations
 randomize phase without directly removing energy. The assumed
 $T_{2q}=30\,\mu\mathrm{s}$ together with $T_{1q}$ gives
 $\gamma_\phi=1/T_{2q}-1/(2T_{1q})$ and
 $L=\sqrt{\gamma_\phi/2}\,\sigma^z$.
 \item \emph{Optional cavity dephasing:} cavity frequency fluctuations
 can be modeled by $L=\sqrt{\eta_j}\,n_j$; $\eta_j=0$ in this baseline.
 Thermal excitation and correlated noise are not included.
\end{itemize}
These are effective Markovian descriptions, not a determination of the
microscopic loss source from the lifetimes. Only the cavity lifetime was
user supplied. Trace out the transmons at the end, keeping leakage
rather than postselecting it away. Energy and correlations follow from
$\langle O\rangle=\tr(\rho O)$ with the appropriate encoding and readout
operators. In particular,
$E=\sum_{i<j}J_{ij}\langle Z_iZ_j\rangle-B\sum_i\langle X_i\rangle$;
the project's spin-glass indicator is
$\chi_{\rm SG}=\frac16\sum_{i,j}\langle Z_iZ_j\rangle^2$
(including $i=j$), computed from the ensemble correlations. It is
neither an extra Hamiltonian term nor a single-copy linear observable.
Quantum-jump sampling provides a memory-efficient alternative to density
matrix integration of the same generator; finite-shot measurement noise
is an additional sampling step, distinct from these physical channels.
\bibliographystyle{apsrev4-2}
\bibliography{references}
\end{document}
